% =========================================================
% Section: Borel Sets
% Chapter: Algebras of Sets — Volume I
% =========================================================
%
% STATUS: Partial.  Definitions and cardinality-independent
% structural facts are complete.  Proofs depending on the
% topology of R (open sets, Gd/Fs hierarchy, Cantor set) are
% stubbed with forward references to the metric topology chapter.
% The non-surjectivity proof (B(R) neq P(R)) is stubbed pending
% the cardinality chapter.
%
% =========================================================

\begin{tcolorbox}[colback=gray!6, colframe=gray!40, arc=2pt,
  left=6pt, right=6pt, top=4pt, bottom=4pt,
  title={\small\textbf{Section Toolkit --- Quick Reference}},
  fonttitle=\small\bfseries]
\small
\begin{tabular}{l l l l}
\toprule
\textbf{Label} & \textbf{Statement} & \textbf{Proof method} & \textbf{Detail} \\
\midrule
D\ref{def:borel-sigma-algebra}
  & $\mathfrak{B}(X) := \sigma(\mathcal{T})$ for topology $\mathcal{T}$
  & Definition & \hyperref[def:borel-sigma-algebra]{$\downarrow$} \\
D\ref{def:gd-fs-sets}
  & $G_\delta$ and $F_\sigma$ sets defined
  & Definition & \hyperref[def:gd-fs-sets]{$\downarrow$} \\
P\ref{prop:borel-closed-gd-fs}
  & Every $G_\delta$ and $F_\sigma$ set is Borel
  & Closure under countable $\cap$ and $\cup$
  & \hyperref[prf:borel-closed-gd-fs]{$\downarrow$} \\
P\ref{prop:borel-not-all-sets}
  & $\mathfrak{B}(\mathbb{R}) \subsetneq \mathcal{P}(\mathbb{R})$
  & \textsc{Stub} --- requires cardinality chapter
  & \hyperref[prf:borel-not-all-sets]{$\downarrow$} \\
\midrule
\multicolumn{4}{l}{\textit{Deferred to metric topology chapter:}} \\
\multicolumn{4}{l}{%
  $\mathfrak{B}(\mathbb{R}) = \sigma(\text{open intervals})
  = \sigma(\text{closed intervals}) = \sigma(\text{half-open intervals})$} \\
\multicolumn{4}{l}{%
  Every open set in $\mathbb{R}$ is a countable union of open intervals} \\
\multicolumn{4}{l}{%
  The Cantor set is Borel (closed set construction)} \\
\bottomrule
\end{tabular}
\end{tcolorbox}

% ---------------------------------------------------------
% Definitions
% ---------------------------------------------------------

\begin{tcolorbox}[colback=propbox, colframe=propborder, arc=2pt,
  left=6pt, right=6pt, top=4pt, bottom=4pt,
  title={\small\textbf{Definition (Borel $\sigma$-Algebra)}},
  fonttitle=\small\bfseries]
\label{def:borel-sigma-algebra}
Let $(X, \mathcal{T})$ be a topological space.  The
\textbf{Borel $\sigma$-algebra} on $X$, denoted $\mathfrak{B}(X)$,
is the $\sigma$-algebra generated by the open sets:
\[
  \mathfrak{B}(X) := \sigma(\mathcal{T}).
\]
Elements of $\mathfrak{B}(X)$ are called \textbf{Borel sets}.
\end{tcolorbox}

\begin{remark*}[Fully quantified form]
$\forall$ topological space $(X,\mathcal{T})$:
$\mathfrak{B}(X) := \sigma(\mathcal{T})$
is the smallest $\sigma$-algebra on $X$ containing every open set.
\end{remark*}

\begin{remark*}[Flash]
\FlashQ{What is the Borel $\sigma$-algebra $\mathfrak{B}(X)$?}
\FlashA{The $\sigma$-algebra generated by the open sets of $X$:
  $\mathfrak{B}(X) = \sigma(\mathcal{T})$.
  It is the smallest $\sigma$-algebra containing all open sets.}
\FlashHint{The definition makes sense for any topological space.
  On $\mathbb{R}$, it is generated by open intervals --- but proving
  that requires knowing what open sets in $\mathbb{R}$ look like
  (metric topology chapter).}
\end{remark*}

\begin{remark*}[Forward reference]
The identification $\mathfrak{B}(\mathbb{R}) = \sigma(\{(a,b) : a<b\})
= \sigma(\{[a,b] : a\leq b\}) = \sigma(\{(a,b] : a<b\})$
is proved in the metric topology chapter once open sets in $\mathbb{R}$
are available as objects.
\end{remark*}

\begin{tcolorbox}[colback=propbox, colframe=propborder, arc=2pt,
  left=6pt, right=6pt, top=4pt, bottom=4pt,
  title={\small\textbf{Definition ($G_\delta$ and $F_\sigma$ Sets)}},
  fonttitle=\small\bfseries]
\label{def:gd-fs-sets}
Let $(X, \mathcal{T})$ be a topological space.
\begin{itemize}
  \item A \textbf{$G_\delta$ set} is a countable intersection of open sets.
  \item An \textbf{$F_\sigma$ set} is a countable union of closed sets.
\end{itemize}
\end{tcolorbox}

\begin{remark*}[Fully quantified form]
$A \subseteq X$ is $G_\delta$ iff $A = \bigcap_{n=1}^\infty U_n$
for some open sets $U_n\in\mathcal{T}$.
$A$ is $F_\sigma$ iff $A = \bigcup_{n=1}^\infty F_n$
for some closed sets $F_n = U_n^c$, $U_n\in\mathcal{T}$.
\end{remark*}

\begin{remark*}[Flash]
\FlashQ{What are $G_\delta$ and $F_\sigma$ sets, and why are they Borel?}
\FlashA{$G_\delta$: countable intersection of opens.
  $F_\sigma$: countable union of closeds.
  Both are Borel: opens are Borel by definition; closed sets are complements
  of opens, hence Borel; $\sigma$-algebras are closed under countable $\cap$
  and $\cup$.}
\FlashHint{The notations come from German: $G$ = Gebiet (region/open set),
  $F$ = Ferm\'e (closed), $\delta$ = Durchschnitt (intersection),
  $\sigma$ = Summe (union).}
\end{remark*}

% ---------------------------------------------------------
% Propositions
% ---------------------------------------------------------

\begin{proposition}[\texorpdfstring{%
  \hyperref[prf:borel-closed-gd-fs]{%
    $G_\delta$ and $F_\sigma$ Sets are Borel}}{%
  Gdelta and Fsigma Sets are Borel}]
\label{prop:borel-closed-gd-fs}
Every $G_\delta$ set and every $F_\sigma$ set in a topological space
$(X,\mathcal{T})$ is a Borel set.
\end{proposition}

\begin{remark*}[Fully quantified form]
$\forall$ topological space $(X,\mathcal{T})$:
if $A$ is $G_\delta$ or $F_\sigma$ then $A\in\mathfrak{B}(X)$.
\end{remark*}

\begin{remark*}[Flash]
\FlashQ{Prove that every $G_\delta$ and $F_\sigma$ set is Borel.}
\FlashA{$G_\delta$: $A = \bigcap_{n=1}^\infty U_n$ with $U_n$ open, so
  $U_n\in\mathfrak{B}(X)$ for all $n$; $\mathfrak{B}(X)$ is closed under
  countable intersection (De Morgan + countable union), so $A\in\mathfrak{B}(X)$.
  $F_\sigma$: $A = \bigcup_{n=1}^\infty F_n$ with $F_n$ closed; each
  $F_n = U_n^c\in\mathfrak{B}(X)$; countable union gives $A\in\mathfrak{B}(X)$.}
\FlashHint{Both cases reduce to closure properties already proved for
  $\sigma$-algebras in Section~4.  No new set machinery is needed.}
\FlashImplies{prop:borel-not-all-sets}
\end{remark*}

% ---------------------------------------------------------
% Stub: B(R) neq P(R)
% ---------------------------------------------------------

\begin{proposition}[\texorpdfstring{%
  \hyperref[prf:borel-not-all-sets]{%
    The Borel Sets Do Not Exhaust $\mathcal{P}(\mathbb{R})$}}{%
  Borel Sets Do Not Exhaust the Power Set}]
\label{prop:borel-not-all-sets}
$\mathfrak{B}(\mathbb{R}) \subsetneq \mathcal{P}(\mathbb{R})$:
there exist subsets of $\mathbb{R}$ that are not Borel sets.
\end{proposition}

\begin{remark*}[Fully quantified form]
$|\mathfrak{B}(\mathbb{R})| = \mathfrak{c}$
whereas $|\mathcal{P}(\mathbb{R})| = 2^{\mathfrak{c}}$,
and $\mathfrak{c} < 2^{\mathfrak{c}}$ by Cantor's theorem.
Therefore $\mathfrak{B}(\mathbb{R})\neq\mathcal{P}(\mathbb{R})$.
\end{remark*}

\begin{remark*}[Forward reference]
\textsc{Proof deferred.}  This result requires:
(i) $|\mathfrak{B}(\mathbb{R})| = \mathfrak{c}$
(a transfinite induction over the Borel hierarchy, or a direct
cardinality count of generated sets),
(ii) $|\mathcal{P}(\mathbb{R})| = 2^{\mathfrak{c}}$ (power set cardinality),
(iii) $\mathfrak{c} < 2^{\mathfrak{c}}$ (Cantor's theorem).
All three ingredients are developed in the cardinality chapter.
The proof is completed there as the capstone result.
\end{remark*}

\begin{remark*}[Flash]
\FlashQ{Why does $\mathfrak{B}(\mathbb{R})\subsetneq\mathcal{P}(\mathbb{R})$?}
\FlashA{Cardinality: $|\mathfrak{B}(\mathbb{R})| = \mathfrak{c}$ but
  $|\mathcal{P}(\mathbb{R})| = 2^{\mathfrak{c}} > \mathfrak{c}$
  by Cantor's theorem.
  So there are strictly more subsets of $\mathbb{R}$ than Borel sets.}
\FlashHint{The proof is not constructive --- it does not exhibit a specific
  non-Borel set.  For an explicit example, the analytic sets provide one,
  but that is well beyond this chapter.}
\FlashImplies{prop:generated-sigma-exists}
\end{remark*}
